\documentclass[11pt]{article}
\usepackage[a4paper,margin=25mm]{geometry}
\usepackage{amsmath,amssymb,booktabs,graphicx,microtype,xcolor}
\usepackage[hidelinks]{hyperref}
\usepackage{enumitem}
\setlist{nosep}
\graphicspath{{../../figs/}}
% Generated values are replaced by docs/numbers.py after cloud integration.
\newcommand{\TotalRuns}{0}
\newcommand{\CompletedRuns}{0}
\newcommand{\NonCompletedRuns}{0}
\newcommand{\SourceCommitCount}{0}


\newcommand{\method}[1]{\textsc{#1}}
\newcommand{\missingfigure}[2]{%
  \begin{figure}[p]
  \centering
  \IfFileExists{../../figs/#1.pdf}{\includegraphics[width=0.94\linewidth]{../../figs/#1.pdf}}{%
    \fbox{\parbox[c][105mm][c]{0.86\linewidth}{\centering\Large #2\\[6mm]\normalsize Cloud evidence has not yet been returned.}}}
  \caption{#2. Missing panels remain explicit rather than being inferred.}
  \end{figure}
  \clearpage}

\title{Deep BSDE and Tensor-Train Methods for High-Dimensional Parabolic PDEs\\
\large Reproducible benchmark and USTC 107 execution protocol}
\author{Quentinluoqz}
\date{Draft generated \today}

\begin{document}
\maketitle

\begin{abstract}
Classical grid discretizations scale exponentially with spatial dimension. This project
compares two alternatives on a common suite of problems with dimensions from 5 to 100:
Deep BSDE solvers, which exploit stochastic representations and neural approximation,
and tensor-train methods, which exploit low-rank structure. The experimental protocol
jointly measures solution error, uncertainty, wall-clock cost, memory, numerical precision,
and hardware throughput. At the current local-development gate, all scientific interfaces,
reference solvers, resumable training, tensor-train representation experiments, and Slurm
orchestration are implemented. Quantitative conclusions are intentionally deferred until
the USTC 107 campaign artifacts have been returned and provenance-validated.
\end{abstract}

\section{Motivation and research questions}
For a tensor-product grid with $n$ points per coordinate, a direct representation requires
$n^d$ values. At $d=100$ this is intractable even for very small $n$. Deep BSDE methods
replace the grid by Monte Carlo paths and trainable gradient approximations. Tensor trains
replace a dense tensor by a chain of low-rank cores. Both can avoid an explicit exponential
representation, but they make different structural and optimization assumptions.

The benchmark addresses four primary questions: how error scales with dimension; how the
cost--accuracy Pareto frontier changes; when low TT rank is retained; and when fp32 or TF32
is sufficient relative to fp64. Seed variance and failure modes are retained as scientific
evidence rather than silently removed.

\section{Problem class and stochastic representation}
We consider semilinear terminal-value problems
\begin{equation}
  \partial_t u + \mu\!\cdot\!\nabla u
  + \tfrac12\operatorname{Tr}(\sigma\sigma^\top\nabla^2u)
  + f(t,x,u,\sigma^\top\nabla u)=0,\qquad u(T,x)=g(x).
\end{equation}
With $X$ following the forward SDE induced by $\mu$ and $\sigma$, the associated backward
equation evolves $Y_t=u(t,X_t)$ and $Z_t=\sigma^\top\nabla u(t,X_t)$. Deep BSDE optimizes
the terminal mismatch $\mathbb{E}|Y_T-g(X_T)|^2$.

\subsection{Benchmark equations}
\begin{table}[ht]
\centering
\small
\begin{tabular}{llll}
\toprule
ID & Type & Reference & Default horizon \\
\midrule
E1 & Black--Scholes--Barenblatt & analytic quadratic solution & $T=1$ \\
E2 & linear-quadratic control & scalar Riccati solution & $T=1$ \\
E3 & quadratic-gradient HJB & Cole--Hopf Monte Carlo, 95\% CI & $T=1$ \\
E4 & Allen--Cahn & literature value at $d=100,x=0$ only & $T=0.3$ \\
E5 & Gaussian heat equation & analytic heat kernel & $T=1$ \\
\bottomrule
\end{tabular}
\caption{The reference boundary is part of the data model. E4 errors outside the canonical
point are undefined, not estimated.}
\end{table}

For E3 the Cole--Hopf transform gives
\begin{equation}
 u(t,x)=-\lambda^{-1}\log\mathbb{E}\left[
 \exp\{-\lambda g(x+\sqrt{2}W_{T-t})\}\right].
\end{equation}
The implementation samples in bounded-memory batches and propagates the standard error
through the logarithm to form a 95\% confidence interval. The production reference uses
$10^7$ samples. Tests compare the one-dimensional case with deterministic quadrature and
verify the expected square-root sample-size behavior of the interval width.

\section{Deep BSDE methods}
\method{M1a} learns $u(0,x_0)$, an initial $Z_0$, and an independent network for every
internal time step. Each step network contains two linear--batch-normalization--ReLU blocks.
\method{M1b} shares a tanh network with input $(t,x)$ across all steps. Both use the same
Euler--Maruyama paths, optimizer lifecycle, result schema, and failure policy.

Checkpoints contain the model, optimizer, scheduler, iteration, loss history, complete
configuration hash, and Python, NumPy, Torch, and CUDA random-number states. Resume rejects
a mismatched configuration. Non-finite losses, gradients, or parameters terminate with a
\texttt{diverged} metric rather than disappearing from the campaign.

The precision modes are deliberately narrow: fp32 uses float32; TF32 keeps float32 tensors
and enables CUDA tensor-core matrix multiplication; fp64 converts parameters, paths, and
equation operations to double. AMP and half precision are outside this comparison.

\section{Tensor-train methods}
Tier 1 measures the compressibility of known solution functions. A rank-adaptive TT-cross
algorithm samples a 32-point Chebyshev grid per dimension without materializing the full
tensor. It uses held-out validation and test indices, records the achieved rank spectrum,
compression, function calls, wall-clock time, and both discrete and continuous interpolation
errors. Failure to reach a target tolerance is recorded as \texttt{not\_reached}.

Tier 2 is a conditional external integration. The Richter--Sallandt--N\"usken solver relies
on xerus/SALSA, which has compiled dependencies and AGPL obligations. No external source is
vendored. A pinned adapter is enabled only after a recorded license and no-sudo build review;
otherwise Tier 1 remains the reproducible tensor baseline.

\section{Experimental design}
The mandatory campaign contains 120 Deep BSDE runs, 45 rank experiments, and 36 precision
runs. Optional accepted Tier-2 integration adds 30 runs. The strong-scaling experiment uses
a fixed global batch for one versus two A100 GPUs. Manifests are immutable, one row is one
atomic command, and each row binds method, equation, dimension, seed, precision, source
commit, and config hash.

\begin{table}[ht]
\centering
\begin{tabular}{lrrl}
\toprule
Campaign & Planned runs & Primary resource & Gate \\
\midrule
Deep BSDE main & 120 & RTX 5090 & pilot E1 then E3 \\
TT rank & 45 & CPU-6530 & E5 low-rank smoke \\
Precision & 36 & RTX 5090 / A100 & matched configs \\
Strong scaling & 6 & A100 & fixed global batch \\
TT Tier 2 & 30 & CPU & license and build review \\
\bottomrule
\end{tabular}
\caption{Planned campaign sizes. The manifest generator is the executable source of truth.}
\end{table}

\section{Platform execution and evidence}
Static partition and QoS values from the training material are treated as defaults only.
Every campaign starts with live \texttt{sinfo}, \texttt{scontrol}, \texttt{sacctmgr}, module,
Python, Torch, CUDA, and storage snapshots. Login nodes perform only preparation and job
management. GPU tests execute inside allocated jobs.

RTX 5090 arrays default to one GPU and four CPU cores with concurrency four; A100 arrays
default to concurrency two. Templates use one task and explicit CPUs per task. Checkpointing,
signal handling, and requeue preserve interrupted work. \texttt{sacct} rows, stdout/stderr,
Grafana screenshots, and environment snapshots are returned with the atomic metrics.
Slurm JWTs remain user-controlled environment variables and are never written to source,
logs, bundles, or chat.

\section{Results protocol}
The local repository currently contains \TotalRuns\ integrated result rows,
\CompletedRuns\ completed rows, and \NonCompletedRuns\ non-completed rows from
\SourceCommitCount\ source commits. These zero-valued draft macros are replaced directly
from the validated summary table after cloud return; they are not hand-edited.

\clearpage
\missingfigure{f1_error_vs_dim}{Relative error versus dimension}
\missingfigure{f2_cost_error_pareto}{Cost--error Pareto frontier}
\missingfigure{f3_tt_rank}{Tensor-train rank versus dimension and tolerance}
\missingfigure{f4_precision_throughput}{Precision, accuracy, and throughput}
\missingfigure{f5_seed_variance}{Deep BSDE seed variance}
\missingfigure{f6_strong_scaling}{One- versus two-GPU strong scaling}

No numerical conclusion is stated until each plotted point has a valid reference boundary,
source commit, configuration hash, and accounting record. Incomplete evidence produces a
visible missing panel or warning instead of an interpolated claim.

\section{Validation and reproducibility}
Fast tests cover configuration rejection, registry behavior, analytic references, batched
tensor shapes in dimensions 1, 5, and 50, Monte Carlo confidence intervals, checkpoint
continuity, manifest expansion, array indexing, and graceful incomplete plotting. GPU and
long numerical tests are explicit markers. A clean clone must pass format, lint, type, and
CPU test gates before a cloud bundle can be created.

\section{Limitations and planned completion}
The current report is a local-stage artifact. It does not claim the target E1 or E3 cloud
accuracy, hardware throughput, scaling efficiency, or a successful Tier-2 build. Those
claims require the user-run USTC campaign and returned evidence. Hyperparameter conclusions
will be restricted to the frozen campaign rather than generalized beyond the tested range.

\section{Conclusion}
The project provides a common scientific and engineering boundary for two distinct ways of
addressing high dimensionality. Its principal design choice is that numerical uncertainty,
failure, resource use, and provenance are first-class outputs. The final comparison will be
written only after the cloud artifacts pass the same boundary.

\appendix
\section{Atomic command contract}
\begin{verbatim}
python scripts/run_one.py --method M --equation E --dim D \
  --seed S --precision P [--resume] [--override key=value]
\end{verbatim}

\section{Required repository gates}
\begin{verbatim}
uv run ruff format --check src tests experiments scripts
uv run ruff check src tests experiments scripts
uv run mypy src
uv run pytest -m "not slow and not gpu"
uv run pre-commit run --all-files
\end{verbatim}

\begin{thebibliography}{9}
\bibitem{e2017} W. E, J. Han, and A. Jentzen. Deep learning-based numerical methods
for high-dimensional parabolic PDEs and BSDEs. \emph{Communications in Mathematics
and Statistics}, 5:349--380, 2017.
\bibitem{han2018} J. Han, A. Jentzen, and W. E. Solving high-dimensional partial
differential equations using deep learning. \emph{PNAS}, 115:8505--8510, 2018.
\bibitem{richter2021} L. Richter, L. Sallandt, and N. N\"usken. Solving high-dimensional
parabolic PDEs using the tensor train format. \emph{PMLR}, 139:8998--9009, 2021.
\end{thebibliography}
\end{document}
